% P15-U3/U6: Kass--Raftery model comparison, with exact finite/dominated proofs.
\chapter{Bayesian Model Comparison}
\label{sec:p15-bayesian-model-comparison}

A Bayes factor compares the probability laws that two models assign to the
observed data after integrating their parameter priors. It is distinct from
a maximized likelihood ratio, a posterior model probability and a
significance-test tail probability. Each answers a different question.

The source for this component is Kass and Raftery, Sections 3, 5, 6.1 and
8.2 \cite[pp.~776--777, 781--784, 789--790]{KassRaftery1995BayesFactors}.
The single-model posterior construction, conjugacy, prediction and stopping
mechanisms belong to \cref{sec:p15-bayesian-updating-prediction}.
Here the parameter priors, observation space and model probabilities are
explicit parts of a comparison. The exact results below do not invoke a
large-sample approximation or silently replace an integral by its maximum.

\section{Models, marginal likelihoods and posterior odds}
\label{sec:p15-model-likelihoods}

\begin{hbdefinition}{def:p15-model-comparison-experiment}{A finite family of dominated models}
Let \(\mathcal K=\{0,\ldots,K\}\), \(K\geq1\), index the model labels.
Let \((\mathcal X,\mathcal A)\) and each \((\Theta_k,\mathcal T_k)\) be
nonempty standard Borel spaces. Fix one sigma-finite observation measure
\(\nu\) on \(\mathcal X\). Model \(H_k\) supplies a proper parameter
prior \(\Pi_k\) and a finite, jointly measurable nonnegative likelihood density
\(f_k(x\mid\theta)\), satisfying
\[
 \int_{\mathcal X}f_k(x\mid\theta)\,\nu(dx)=1
 \qquad(\theta\in\Theta_k).
\]
A likelihood is evaluated as a function of \(\theta\) after observing \(x\);
the same function is a density in \(x\) before observation.

The \emph{model marginal likelihood}, or prior predictive density, is
\[
 m_k(x)=\int_{\Theta_k}f_k(x\mid\theta)\,\Pi_k(d\theta).
\]
Write \(Q_k(A)=\int_A m_k\,d\nu\). Model prior probabilities
\(\alpha_k\geq0\) satisfy \(\sum_k\alpha_k=1\). Their predictive mixture
is \(Q_\alpha=\sum_k\alpha_kQ_k\), with density
\(e(x)=\sum_{k:\alpha_k>0}\alpha_km_k(x)\); zero prior weights do not
require forming \(0\cdot\infty\).
Use the comparison domain
\[
 \mathcal G=\{x:m_k(x)<\infty\text{ for all }k,\quad 0<e(x)<\infty\}.
\]
On \(\mathcal G\), define
\[
 w_k(x)=\frac{\alpha_km_k(x)}{e(x)},\qquad
 B_{ij}(x)=\frac{m_i(x)}{m_j(x)}\quad\text{when }m_j(x)>0.
\]
The first index in \(B_{ij}\) identifies the numerator model, as in the
source. A zero numerator gives factor zero. A zero denominator is not
assigned an ordinary finite factor, and \(0/0\) is left undefined.
\end{hbdefinition}

\begin{hbtheorem}{thm:p15-model-posterior-odds}{proved=annex:p15.model-posterior-odds}{Marginal laws and model odds}
Each \(m_k\) is measurable, integrates to one and is finite \(\nu\)-almost
everywhere; \(Q_\alpha(\mathcal G)=1\). Extending \(w_k=\alpha_k\)
outside \(\mathcal G\) gives a version of the conditional probabilities
of the model label given the data. For \(x\in\mathcal G\), where
\(\alpha_j>0\) and \(m_j(x)>0\),
\[
 \frac{w_i(x)}{w_j(x)}
 =B_{ij}(x)\frac{\alpha_i}{\alpha_j}.
\]
On \(\mathcal G\), where both marginal likelihoods are positive,
\(B_{ij}B_{ji}=1\) and \(\log B_{ij}=\log m_i-\log m_j\).
Bayes factors hold the within-model priors fixed and do not involve
\(\alpha\); posterior model probabilities do involve \(\alpha\).

If \(x\in\mathcal G\) and \(m_k(x)>0\), the parameter posterior within \(H_k\) is the kernel
from \cref{thm:p15-bayes-posterior-kernel}. The joint posterior of model
and parameter weights that kernel by \(w_k(x)\). A prior-zero model has
posterior probability zero on \(\mathcal G\).
\end{hbtheorem}

For two models with positive prior probabilities, putting \(H_0\) in
the numerator gives
\[
 w_0(x)=\frac{\alpha_0 B_{01}(x)}
              {\alpha_0 B_{01}(x)+\alpha_1}
 \qquad(m_0(x),m_1(x)>0).
\]
Only for equal model priors does the Bayes factor equal posterior odds.
Even then it is an odds ratio, not a probability. A factor of 20 with
prior odds \(1:100\) gives posterior odds \(1:5\), or probability \(1/6\)
for the numerator model.

The logarithmic identity also compares the predictive log densities assigned
by the two models; it remains a numerical comparison when both are regarded
as approximations. Its sequential predictive interpretation uses the
factorization in \cref{sec:p15-bayesian-updating-prediction}.
No nesting condition on the parameter spaces is needed for a Bayes factor.

Kass and Raftery's interpretive labels are rough conventions, not calibrated
posterior probabilities or decision rules. Their combined Jeffreys scale
uses approximate factor ranges 1--3.2, 3.2--10, 10--100 and above 100,
with descriptions ranging from negligible mention to substantial, strong
and decisive evidence. Their modified \(2\log B_{10}\) scale uses 0--2,
2--6, 6--10 and above 10, corresponding approximately to factors 1--3,
3--20, 20--150 and above 150. The latter descriptions are negligible
mention, positive, strong and very strong. The numerical boundaries are
rounded and their interpretation is contextual
\cite[Sec.~3.2]{KassRaftery1995BayesFactors}.

\section{Which changes preserve a comparison?}
\label{sec:p15-model-density-invariance}

\begin{hbtheorem}{thm:p15-model-density-invariance}{proved=annex:p15.model-density-invariance}{Density versions, reference measures and parameter coordinates}
Changing jointly measurable likelihood versions without changing their
conditional data laws preserves marginal laws and marginal densities
\(\nu\)-almost everywhere. The resulting posterior model probabilities
agree \(Q_\alpha\)-almost everywhere.

More intrinsically, put \(Q_*=(K+1)^{-1}\sum_kQ_k\), and let
\(r_k=dQ_k/dQ_*\). Then
\[
 B_{ij}=\frac{r_i}{r_j}\quad\text{where }r_j>0,
 \qquad
 w_k=\frac{\alpha_kr_k}{\sum_\ell\alpha_\ell r_\ell}
\]
on the appropriate almost-everywhere domains. Thus any common sigma-finite
observation measure dominating the same model marginal laws gives the same
comparison \(Q_*\)-almost everywhere, and the same posterior
\(Q_\alpha\)-almost everywhere.

For the particular change \(d\nu'=h\,d\nu\), where \(0<h<\infty\)
almost everywhere and \(\nu'\) is sigma-finite, compatible versions are
\(f'_k=f_k/h\) and \(m'_k=m_k/h\). Their factors agree pointwise where
these equalities and positive finite denominators hold.

A bimeasurable bijection of a parameter space also preserves \(m_k(x)\)
pointwise when both its prior and its likelihood are transported by that
bijection. Declaring a new flat prior in the new coordinates is not such
a transport in general.

A bimeasurable bijection \(S:\mathcal X\to\mathcal X'\) of observation
spaces has pushed-forward reference measure
\((S_\#\nu)(A)=\nu(S^{-1}A)\). It preserves the factors on corresponding
data when \(\nu'=S_\#\nu\)
and \(f'_k(y\mid\theta)=f_k(S^{-1}y\mid\theta)\). In that case
\(m'_k(Sx)=m_k(x)\). Using another dominating measure afterward invokes
the reference-measure qualification above.
\end{hbtheorem}

Continuous data require this version qualification. For example, two models
can both have the uniform law on \([0,1]\), with density one. Altering the
second density only at \(x=1/2\) to an arbitrary finite \(c>0\) leaves
its probability law unchanged but changes the pointwise formula for
\(B_{10}(1/2)\) to \(c\). The almost-everywhere factor is still one.
Accordingly, the normal calculations below specify the usual continuous
density versions; they do not claim that arbitrary modifications at a null
point have an invariant numerical value.

The data must also mean the same thing in each model. A density for an
ordered sample, a mass for a count and a likelihood for selected observations
are different observation descriptions. A comparison can use one of them
consistently or prove an appropriate common-factor reduction; it cannot
silently mix their reference measures.

\section{Integrating nuisance parameters instead of maximizing}
\label{sec:p15-nuisance-model-integration}

\begin{hbtheorem}{thm:p15-nuisance-integration}{proved=annex:p15.nuisance-integration}{Averaging a nuisance prior}
Suppose \(\Theta_k=\Psi_k\times\mathcal N_k\) for nonempty standard
Borel spaces, \(\theta=(\psi,\lambda)\), and the declared proper prior has
the form
\(\Pi_k(d\psi,d\lambda)=\Lambda_k(\psi,d\lambda)\Pi_k^\psi(d\psi)\),
where \(\Lambda_k\) is a probability kernel. For the nonnegative measurable
likelihood define
\[
 \bar f_k(x\mid\psi)=\int f_k(x\mid\psi,\lambda)\Lambda_k(\psi,d\lambda).
\]
Then
\[
 m_k(x)=\int\bar f_k(x\mid\psi)\Pi_k^\psi(d\psi).
\]
The inner integral is a prior-weighted average; it is bounded above by
\(\sup_\lambda f_k(x\mid\psi,\lambda)\) whenever that supremum is finite.
Replacing the integral by the supremum is a different operation. Even if
the supremum is measurable, it need not be a normalized data density.
No equality of nuisance priors in different models is automatic.
\end{hbtheorem}

For a calculated example, let \(X\in\{0,1\}\). Model \(H_0\) assigns
\(\Pr(X=1)=.6\). Under \(H_1\), a nuisance label has equal prior weights
and gives success probabilities .9 and .1. Hence \(m_1(1)=m_1(0)=.5\);
at \(X=1\), \(B_{10}=5/6\). Maximizing over the nuisance label gives
.9 at both outcomes, whose sum is 1.8. Its ratio to \(m_0(1)\) is
\(3/2\), reversing the comparison, but it is not the Bayes factor.
If the nuisance prior weight on .9 is \(\beta\), its actual marginal
success probability is \(.1+.8\beta\). This identifies precisely which
prior choice changes the integrated evidence.

The source discusses both matching a shared nuisance prior to the
alternative's marginal and matching it to a conditional prior on a null
constraint. These are specifications to justify, not identities required
by Bayes' theorem \cite[Sec.~5.1, eq.~(13)]{KassRaftery1995BayesFactors}.
A constraint of prior probability zero does not by itself determine which
conditional-prior version to use there; an explicit example appears below.

\section{A normal point null and a proper normal alternative}
\label{sec:p15-normal-bayes-factor}

Let \(\theta_0,\mu_1\in\mathbb R\). Let one observed summary \(X\),
conditional on a real parameter \(\theta\), have law
\(\operatorname{normal}(\theta,s^2)\), where \(s>0\) is known.
Under \(H_0\), \(\theta=\theta_0\). Under \(H_1\),
\(\theta\sim\operatorname{normal}(\mu_1,\tau^2)\), with \(\tau>0\).
Normal notation and its Lebesgue density are those of
\cref{def:p15-normal-law}. Both within-model priors are proper, including
the point mass under \(H_0\).

\begin{hbtheorem}{thm:p15-normal-point-null-factor}{proved=annex:p15.normal-point-null-factor}{The exact normal factor}
With the continuous density versions just specified,
\[
 m_0(x)=\phi_{\theta_0,s^2}(x),\qquad
 m_1(x)=\phi_{\mu_1,s^2+\tau^2}(x),
\]
and
\[
 B_{01}(x)=
 \sqrt{\frac{s^2+\tau^2}{s^2}}\,
 \exp\!\left\{-\frac{(x-\theta_0)^2}{2s^2}
             +\frac{(x-\mu_1)^2}{2(s^2+\tau^2)}\right\}.
\]
For a prior centered at the null, \(\mu_1=\theta_0\), put
\(z=(x-\theta_0)/s\) and \(r=\tau^2/s^2>0\). Then
\[
 B_{01}(x)=\sqrt{1+r}\exp\!\left(-\frac{rz^2}{2(1+r)}\right).
\]
The full sample \(X_1,\ldots,X_n\), \(n\geq1\), conditionally independent
\(\operatorname{normal}(\theta,\sigma^2)\) with known \(\sigma>0\),
gives this same factor at \(x=\bar X\), \(s^2=\sigma^2/n\), because
its remaining likelihood factor is common to both models.
\end{hbtheorem}

A point null is a model with positive prior model mass, if it is to acquire
positive posterior model mass. It is different from the event
\(\theta=\theta_0\) inside the continuous alternative: that event has prior
and posterior probability zero under \(H_1\). A high density near
\(\theta_0\) is not a point probability.

\begin{hbtheorem}{thm:p15-normal-factor-sensitivity}{proved=annex:p15.normal-factor-sensitivity}{Prior scale, normal-class bounds and limiting paths}
For fixed \(x,s,\theta_0\) and a null-centered alternative,
\(B_{01}\to1\) as \(r\downarrow0\), and \(B_{01}\to\infty\) as
\(r\to\infty\). Over all positive normal prior variances,
\[
 \sup_{r>0} B_{10}=
 \begin{cases}
 1,&z^2\leq1,\\
 \displaystyle\frac{\exp((z^2-1)/2)}{|z|},&z^2>1.
 \end{cases}
\]
For \(z^2>1\) the maximum occurs at \(r=z^2-1\); otherwise the
supremum is approached as \(r\downarrow0\). Mixing these centered normal
priors with any proper mixing probability cannot exceed the same bound.
Over all proper priors on \(\theta\), without the centering restriction,
the supremum of \(B_{10}\) is \(\exp(z^2/2)\), attained by a point prior
at \(x\). Selecting such a prior after observing \(x\) is a sensitivity
bound, not a prespecified Bayesian analysis.

For fixed \(\tau,\sigma>0\), along summaries
\(\bar x_n=\theta_0+z\sigma/\sqrt n\) with \(z\) fixed,
\(B_{01}\to\infty\). Along \(\bar x_n=\theta_0+\delta\) with
fixed \(\delta\ne0\), instead \(B_{01}\to0\).
\end{hbtheorem}

These different limits must not be conflated. Making an alternative prior
diffuse for fixed data spreads its predictive probability over outcomes not
observed, which can favor the null. Increasing sample size while fixing a
nonzero observed departure gives a different comparison. The sensitivity
analysis concerns a declared family of proper priors, not a data-independent
justification for whichever member maximizes the factor
\cite[Secs.~5.1--5.2]{KassRaftery1995BayesFactors}.

\begin{hbexample}{def:p15-normal-factor-table}{The same standardized observation under different priors}
For \(z=2\) and equal model prior probabilities:
\[
\begin{array}{r|rr}
 r& B_{01}&w_0\\\hline
 1&.5203&.3422\\
 3&.4463&.3086\\
 9&.5227&.3433\\
 99&1.3807&.5800\\
 999&4.2882&.8109
\end{array}
\]
Values are rounded evaluations of the exact formulas, not observations
from a study. At \(r=3\), the alternative attains its strongest support
within the centered-normal family, yet the equal-prior posterior null
probability is still about .309.
\end{hbexample}

\section{Posterior model averaging}
\label{sec:p15-posterior-model-averaging}

\begin{hbdefinition}{def:p15-model-average-law}{A common quantity across models}
For each \(k\), let \(R_k(x,\cdot)\) be a probability kernel on the same
measurable outcome space \(\mathcal Z\), describing a posterior quantity
or prediction conditional on \(H_k\). For example, a measurable map
\(T_k:\Theta_k\to\mathcal Z\) can send each parameter posterior to a
common quantity. The physical quantity, units and event meanings must
agree across models. If a parameter is absent from one model, assigning
it value zero is an additional model definition, not an implicit rule.
Set
\[
 R(x,A)=\sum_kw_k(x)R_k(x,A).
\]
On exceptional data outside \(\mathcal G\), use the already specified
default weights and the declared kernels.
\end{hbdefinition}

\begin{hbtheorem}{thm:p15-model-average-moments}{proved=annex:p15.model-average-moments}{Mixture laws and total posterior variance}
The function \(R\) is a probability kernel. For a real common quantity
and a fixed \(x\), suppose every component with \(w_k(x)>0\) has finite
second moment. Let its mean and variance be \(a_k,v_k\), and sum only
over these positive-weight components. Then
\[
 a=\sum_kw_ka_k,\qquad
 \operatorname{Var}_R(Z)=\sum_kw_kv_k+\sum_kw_k(a_k-a)^2.
\]
Equivalently the variance is
\(\sum_kw_k(v_k+a_k^2)-a^2\), as in the source.
For an integrable common utility or other function \(u\),
\(\int u\,dR=\sum_kw_k\int u\,dR_k\).
The between-model term is nonnegative. Omitting it from the weighted
within-model variance strictly understates the mixture variance when
positive-weight component means differ, and otherwise leaves it unchanged.
The variance
of a particular selected model need not always be smaller than the mixture.
\end{hbtheorem}

With equal posterior weights and component laws
\(\operatorname{normal}(-2,1)\), \(\operatorname{normal}(2,1)\), the
mixture mean is zero and variance is \(1+4=5\). Selecting either component
reports variance one. Conversely, equal weights on
\(\operatorname{normal}(0,100)\) and the point mass at zero give mixture
variance 50; selecting the first reports 100. These examples locate the
precise omitted-uncertainty statement without claiming a universal ordering
for every possible selected model
\cite[Sec.~6.1, eqs.~(17)--(19)]{KassRaftery1995BayesFactors}.

Model averaging accounts for uncertainty within the declared collection.
It does not show that the collection contains an adequate description of
the data. Relative model evidence and checks for failures of all the models
address different questions, as the source notes when discussing goodness
of fit \cite[Sec.~8.2]{KassRaftery1995BayesFactors}.

\section{Model priors, duplicated labels and improper conventions}
\label{sec:p15-model-prior-conventions}

\begin{hbtheorem}{thm:p15-model-label-refinement}{proved=annex:p15.model-label-refinement}{Preserving prior mass when a label is split}
Keep the model likelihoods and within-model priors fixed. Replacing
\(\alpha\) by another model prior changes posterior odds by the corresponding
prior-odds multiplier, while leaving Bayes factors unchanged.

Replace a model \(H_j\) by finitely many labels having the same marginal
law and the same common-quantity kernels. If their prior masses sum to
\(\alpha_j\), the aggregate posterior probability of that group and the
model-averaged law are unchanged almost everywhere under the original
predictive mixture. Assigning equal mass afresh to every
expanded label need not preserve the original comparison.
\end{hbtheorem}

For example, two identical predictive laws with equal model masses have
posterior masses \(1/2,1/2\). Copy the first label nine times and assign
uniform mass to the resulting ten labels: its group now has posterior .9.
If its marginal likelihood instead is twice the second's, the original
equal-prior group probability is \(2/3\), but the nine-copy uniform encoding
gives \(18/19\). The changed posterior support came from changing prior
group mass; the Bayes factor did not change. Uniformity over names therefore
requires attention to what counts as a distinct model.

\begin{hbtheorem}{thm:p15-improper-model-priors}{proved=annex:p15.improper-model-priors}{Arbitrary constants and null-constraint priors}
Let \(\widetilde\Pi_k\) be a sigma-finite measure of infinite mass, and
suppose \(0<Z_k(x)=\int f_k(x\mid\theta)\widetilde\Pi_k(d\theta)<\infty\)
at the data being used. Multiplying this measure by \(c_k>0\) leaves
the normalized within-model posterior unchanged, but multiplies \(Z_k\)
by \(c_k\). Consequently \(Z_i/Z_j\) depends on the arbitrary ratio
\(c_i/c_j\). Also \(\int Z_k\,d\nu=\infty\), so \(Z_k\) is not a
proper prior predictive density.

A common multiplicative constant cancels if the same explicitly identified
nuisance measure is used in both integrations. This algebra alone neither
normalizes the predictive measures nor justifies a unique choice of nuisance
measure.

Even proper priors do not automatically specify a nuisance prior on a
prior-null constraint. For independent uniform \(\beta,\psi\in[0,1]\),
both a conditional law of \(\beta\) that is uniform everywhere except
\(\psi=1/2\), where it is \(\delta_0\), and one that is instead
\(\delta_1\) there are valid versions. Using these versions as priors for
a separate null model \(\psi=1/2\) makes different models.
\end{hbtheorem}

For \(X\mid\theta\sim\operatorname{normal}(\theta,1)\) and a formal
flat measure \(c\,d\theta\), the integrated likelihood is \(c\) for every
\(x\), while the normalized parameter posterior is
\(\operatorname{normal}(x,1)\). Comparing it with a point-null predictive
density \(\phi_{0,1}(x)\) produces the arbitrary value
\(c/\phi_{0,1}(x)\). A proper uniform prior on \([-A,A]\), in contrast,
has marginal density at most \(1/(2A)\), tending to zero for fixed data.
It does not provide a distinguished finite value of \(c\).

For the proper null-constraint example, suppose data are binary with
\(\Pr(X=1\mid\beta,\psi)=.2+.6\beta\). The two permitted conditional
versions at \(\psi=1/2\) give null-model success probabilities .2 and .8.
Neither choice follows uniquely from the joint alternative prior.
A substantive compatibility convention can choose one; a probability-zero
conditioning event alone cannot
\cite[Secs.~5.1, 5.3]{KassRaftery1995BayesFactors}.

\section{Finite evaluation and likelihood constants}
\label{sec:p15-finite-model-evaluation}

\begin{hbtheorem}{thm:p15-finite-model-evaluation}{proved=annex:p15.finite-model-evaluation}{Exact sums and permitted cancellations}
If model \(k\) has finite parameter values \(\theta_{k\ell}\) with
prior masses \(p_{k\ell}\), then, at a datum where all the displayed
likelihood values are finite,
\[
 m_k(x)=\sum_\ell p_{k\ell}f_k(x\mid\theta_{k\ell}),\qquad
 w_k(x)=\frac{\alpha_km_k(x)}{\sum_j\alpha_jm_j(x)}
\]
on the positive finite comparison domain. A factor \(c(x)\in(0,\infty)\)
common to all model likelihoods cancels from every Bayes factor and
posterior model weight. Factors depending on the model do not generally
cancel, even if they are constant in that model's parameters.

For finite nonnegative terms \(a_{k\ell}=\alpha_kp_{k\ell}
f_k(x\mid\theta_{k\ell})\), with at least one positive term, set
\(t_{k\ell}=\log a_{k\ell}\) only for positive terms and
\(t_*=\max t_{k\ell}\). Then
\[
 w_k(x)=
 \frac{\sum_{\ell:a_{k\ell}>0}\exp(t_{k\ell}-t_*)}
      {\sum_{j,\ell:a_{j\ell}>0}\exp(t_{j\ell}-t_*)}.
\]
This is an exact real-arithmetic rescaling identity, not a guarantee about
rounding or underflow in an arbitrary implementation.
\end{hbtheorem}

As a finite example, observe one success in two Bernoulli trials as a
binomial count. Under \(H_0\), success probability is .5; under \(H_1\),
it is .2 or .8 with equal prior masses. The full mass values are
\(m_0=2(.5)(.5)=.5\) and
\(m_1=\tfrac12[2(.2)(.8)+2(.8)(.2)]=.32\).
Thus \(B_{10}=.64=16/25\), and equal model priors give \(w_1=16/41\).
Removing the common binomial coefficient two preserves this ratio.

For contrast, at \(x=0\), the continuous densities of
\(\operatorname{normal}(0,1)\) and \(\operatorname{normal}(0,4)\)
have ratio two. Dropping their variance-dependent normalizing constants
leaves two exponential kernels equal to one and incorrectly gives ratio
one. Every calculation must retain or explicitly justify cancellation of
the relevant constants \cite[Sec.~3.1]{KassRaftery1995BayesFactors}.

\section{Significance and posterior probability}
\label{sec:p15-bayes-factor-significance}

In the known-variance normal point-null experiment, let
\(\Phi\) be the standard normal CDF and put
\[
 p_{\rm tail}(x)=2\{1-\Phi(|z|)\},\qquad z=(x-\theta_0)/s.
\]
This is the null tail probability for an observation at least as extreme
in absolute standardized value as the one obtained. Its probability law
is computed under \(H_0\); it is not conditional probability of \(H_0\).

\begin{hblemma}{thm:p15-normal-significance-comparison}{proved=annex:p15.normal-significance-comparison}{A calibrated null-tail statistic need not equal posterior model probability}
Under \(H_0\), for \(0<u<1\),
\[
 \Pr_0\{p_{\rm tail}(X)\leq u\}=u.
\]
At \(z=2\) the value is approximately .0455, unchanged across the
alternative priors in \cref{def:p15-normal-factor-table}, whose posterior
null probabilities differ. Along the fixed-\(z\), increasing-\(n\) path
in \cref{thm:p15-normal-factor-sensitivity}, the tail probability is
constant while \(w_0\to1\) for fixed positive model prior masses.
For a fixed nonzero observed departure, the different limit stated there
applies instead.
\end{hblemma}

The comparison preserves both meanings: the null-tail statistic has its
specified repeated-sampling calibration, while a model posterior incorporates
the alternative's predictive law and the prior model odds. Changing those
priors leaves the tail statistic unchanged. Conversely, a Bayes factor used
as a frequentist test statistic needs a separate repeated-sampling calibration.
Neither its name nor an evidence-label threshold supplies that calibration
\cite[Sec.~8.2]{KassRaftery1995BayesFactors}.

This component supplies exact dominated and finite calculations. The source's
Laplace and Schwarz approximations, general nuisance-orthogonality
asymptotics, generalized-linear-model prior recipes, training-sample methods
for improper priors, and algorithms for very large model spaces require their
own hypotheses and derivations. They are not consequences of the exact
normal formula. Single-model sequential prediction and the conditions for
ignoring a stopping mechanism remain with
\cref{sec:p15-bayesian-updating-prediction}.
